% Section publicaciones
\sectionTitle{Publicaciones}{\faBook}

%%%%%%%%%%%%%%%%%%%%%%%%%%%%%%%%%%%%%%%%
% INSTRUCCIONES DE USO
%%%%%%%%%%%%%%%%%%%%%%%%%%%%%%%%%%%%%%%%
% 
% IMPORTANTE: Para que esta sección funcione correctamente:
% 1. Descomenta \addbibresource{my_publications.bib} en cv.tex
% 2. Crea/edita el archivo my_publications.bib con tus publicaciones
% 3. Compila con: lualatex → biber → lualatex → lualatex
%
% FORMATOS DE COMPILACIÓN:
%   lualatex cv.tex
%   biber cv
%   lualatex cv.tex
%   lualatex cv.tex
%
%%%%%%%%%%%%%%%%%%%%%%%%%%%%%%%%%%%%%%%%

\begin{publications}
	
	%%%%%%%%%%%%%%%%%%%%%%%%%%%%%%%%%%%%%%%%
	% OPCIÓN 1: LISTA COMPLETA SIN CATEGORÍAS (más simple)
	%%%%%%%%%%%%%%%%%%%%%%%%%%%%%%%%%%%%%%%%
	% Usar esta opción si tienes pocas publicaciones (1-5)
	% o si no quieres separar por tipo
	
	% Descomentar la siguiente línea para mostrar todas las publicaciones juntas:
	%\nocite{*}
	%\printbibliography[heading={none}]
	
	%%%%%%%%%%%%%%%%%%%%%%%%%%%%%%%%%%%%%%%%
	% OPCIÓN 2: SEPARADO POR TIPOS DE PUBLICACIÓN (recomendado)
	%%%%%%%%%%%%%%%%%%%%%%%%%%%%%%%%%%%%%%%%
	% Usar esta opción si tienes varias publicaciones de diferentes tipos
	% Permite organizar mejor tu producción académica
	
	\nocite{*} % Incluir todas las publicaciones del .bib sin citarlas en el texto
	
	% Artículos en revistas científicas (journals)
	\printbibliography[heading=subbibliography,title={Artículos de Revistas},type=article]
	
	% Ponencias en conferencias y congresos
	\printbibliography[heading=subbibliography,title={Actas de Conferencias},type=inproceedings]
	
	% Libros y capítulos de libros
	\printbibliography[heading=subbibliography,title={Libros y Capítulos},filter={booksandchapters}]
	
	% Working papers y documentos de trabajo
	%\printbibliography[heading=subbibliography,title={Documentos de Trabajo},type=unpublished]
	
	% Tesis (maestría, doctorado)
	%\printbibliography[heading=subbibliography,title={Tesis},type=thesis]
	
	% Reportes técnicos
	%\printbibliography[heading=subbibliography,title={Reportes Técnicos},type=techreport]
	
	%%%%%%%%%%%%%%%%%%%%%%%%%%%%%%%%%%%%%%%%
	% OPCIÓN 3: PUBLICACIONES SELECCIONADAS (manual)
	%%%%%%%%%%%%%%%%%%%%%%%%%%%%%%%%%%%%%%%%
	% Usar esta opción si quieres mostrar solo algunas publicaciones específicas
	% sin usar archivos .bib
	
	%\begin{itemize}
	%    \item \textbf{Achalma, E.} (2024). ``Determinantes de la informalidad laboral en Ayacucho: Un análisis econométrico 2018-2023''. \emph{Revista de Economía Regional}, 15(2), 45-67. DOI: 10.1234/rer.2024.001
	%    
	%    \item \textbf{Achalma, E.}, \& García, J. (2023). ``Evaluación de impacto del Programa Juntos en la desnutrición infantil''. En \emph{Actas del XIX Congreso de Economistas}, 120-135. Lima: CIES.
	%    
	%    \item \textbf{Achalma, E.} (2023). ``Machine Learning aplicado a la predicción de morosidad crediticia en microfinanzas''. \emph{Documento de Trabajo No. 45}, Centro de Investigación Económica, UNSCH.
	%\end{itemize}
	
	%%%%%%%%%%%%%%%%%%%%%%%%%%%%%%%%%%%%%%%%
	% OPCIÓN 4: SEPARADO POR AÑO (orden cronológico inverso)
	%%%%%%%%%%%%%%%%%%%%%%%%%%%%%%%%%%%%%%%%
	% Útil si tienes muchas publicaciones y quieres ordenar por año
	
	%\nocite{*}
	%\printbibliography[heading=subbibliography,title={2024},year=2024]
	%\printbibliography[heading=subbibliography,title={2023},year=2023]
	%\printbibliography[heading=subbibliography,title={2022},year=2022]
	
	%%%%%%%%%%%%%%%%%%%%%%%%%%%%%%%%%%%%%%%%
	% OPCIÓN 5: SEPARADO POR ÁREA TEMÁTICA
	%%%%%%%%%%%%%%%%%%%%%%%%%%%%%%%%%%%%%%%%
	% Requiere agregar keywords en el archivo .bib
	% Ejemplo en .bib: keywords = {economia-desarrollo}
	
	%\nocite{*}
	%\printbibliography[heading=subbibliography,title={Economía del Desarrollo},keyword={economia-desarrollo}]
	%\printbibliography[heading=subbibliography,title={Econometría Aplicada},keyword={econometria}]
	%\printbibliography[heading=subbibliography,title={Data Science},keyword={data-science}]
	
	%%%%%%%%%%%%%%%%%%%%%%%%%%%%%%%%%%%%%%%%
	% OPCIÓN 6: PUBLICACIONES REVISADAS POR PARES vs NO REVISADAS
	%%%%%%%%%%%%%%%%%%%%%%%%%%%%%%%%%%%%%%%%
	% Útil para distinguir rigor académico
	
	%\nocite{*}
	%\printbibliography[heading=subbibliography,title={Publicaciones Revisadas por Pares},keyword={peer-reviewed}]
	%\printbibliography[heading=subbibliography,title={Otras Publicaciones},notkeyword={peer-reviewed}]
	
\end{publications}

%%%%%%%%%%%%%%%%%%%%%%%%%%%%%%%%%%%%%%%%
% NOTAS ADICIONALES
%%%%%%%%%%%%%%%%%%%%%%%%%%%%%%%%%%%%%%%%
%
% FILTROS PERSONALIZADOS:
% Para crear filtros personalizados (como booksandchapters), 
% agrega en el preámbulo de cv.tex:
%
% \defbibfilter{booksandchapters}{
	%   type=book or
	%   type=inbook or
	%   type=incollection
	% }
%
% ORDENAMIENTO:
% Por defecto las publicaciones se ordenan por año (más reciente primero)
% Puedes cambiar esto en las opciones de biblatex
%
% ESTILOS DE CITACIÓN:
% Puedes cambiar el estilo de las referencias:
% - apa (American Psychological Association)
% - ieee (Institute of Electrical and Electronics Engineers)
% - chicago (Chicago Manual of Style)
% - vancouver (Vancouver system)
%
% Para cambiar el estilo, modifica en cv.tex:
% \usepackage[style=apa]{biblatex}
%
%%%%%%%%%%%%%%%%%%%%%%%%%%%%%%%%%%%%%%%%

%%%%%%%%%%%%%%%%%%%%%%%%%%%%%%%%%%%%%%%%
% TIPS PARA ACADÉMICOS
%%%%%%%%%%%%%%%%%%%%%%%%%%%%%%%%%%%%%%%%
%
% 1. ÍNDICES DE IMPACTO: Menciona si la revista está indexada
%    Ejemplo: "Publicado en revista indexada en Scopus/WoS"
%
% 2. FACTOR DE IMPACTO: Agrega el cuartil si es relevante
%    Ejemplo: "Revista Q1 en Economics (JCR 2023)"
%
% 3. DOI: Siempre incluye el DOI de tus publicaciones
%    Se puede agregar como url en el archivo .bib
%
% 4. ORCID: Asegúrate de tener tu ORCID actualizado con todas tus publicaciones
%
% 5. ACCESO ABIERTO: Si el artículo es open access, menciónalo
%
% 6. PREPRINTS: Puedes incluir preprints (arXiv, SSRN) en "Documentos de Trabajo"
%
%%%%%%%%%%%%%%%%%%%%%%%%%%%%%%%%%%%%%%%%